\customSection[sec:analysis]{Аналіз предметної області}{0em}

\customSubSection[sec:analysis_two_level_section]{Оформлення програмного коду}{0em}

Оформлення коду є важливим аспектом при написанні технічної документації. Нижче наведено приклади реалізації клієнтської та серверної частин системи управління СТО.

% ----------------------
% HTML Structure
% ----------------------
Фрагмент HTML-коду форми для додавання нової заявки адміністратором:
\begin{lstlisting}[style=nocodeframe, language=HTML]
<div class="form-grid">
    <input id="client" type="text" placeholder="ПІБ Клієнта">
    <input id="car" type="text" placeholder="Авто">
    <select id="mechanic">
        <option value="Майстер Іван">Майстер Іван</option>
        <option value="Майстер Володимир">Майстер Володимир</option>
    </select>
    <input id="labor_cost" type="number" placeholder="Робота (грн)">
    <input id="parts_cost" type="number" placeholder="Запчастини (грн)">
    <button class="btn-save" onclick="addOrder()">Зберегти</button>
</div>
\end{lstlisting}

% ----------------------
% Python Backend
% ----------------------
Реалізація моделі даних та API-запиту на мові Python (FastAPI):
\begin{lstlisting}[style=nocodeframe, language=Python]
class Order(Base):
    __tablename__ = "orders"
    id = Column(Integer, primary_key=True, index=True)
    client_name = Column(String)
    car_info = Column(String)
    labor_cost = Column(Float, default=0.0)
    parts_cost = Column(Float, default=0.0)
    status = Column(String, default="В роботі")

@app.post("/add-order")
async def add_order(order_data: OrderCreate, db: Session = Depends(get_db)):
    new_order = Order(**order_data.dict())
    db.add(new_order)
    db.commit()
    return {"status": "success"}
\end{lstlisting}

Система забезпечує автоматизацію роботи адміністратора, дозволяючи вести централізований облік клієнтів та фінансових витрат на ремонт автомобілів.

\customSubSection[sec:analysis_two_level_section]{Оформлення рисунків}{1em}

Для візуалізації роботи системи наведено приклад інтерфейсу панелі керування.

\setfigure[fig:dashboard_view]{0.8\textwidth}{assets/img}{Інтерфейс панелі керування СТО}

На \textbf{рис.~\ref{fig:dashboard_view}} показано головний екран системи, де адміністратор може керувати списком замовлень та змінювати їх статуси.

\customSubSection[sec:analysis_third_level_section]{Оформлення таблиць}{1em}

Для аналізу структури збережених даних нижче наведено опис основних полів бази даних.

\begin{spacing}{1.3}
    \noindent
    \begin{xltabular}{\textwidth}{| c | >{\raggedright\arraybackslash}X | c |}
        \caption{Структура таблиці orders у БД SQLite} \label{tab:db_orders} \\
        \hline
        \textbf{№} & \textbf{Назва поля} & \textbf{Тип даних} \\ \hline
        \endfirsthead

        \multicolumn{3}{r}{\textit{}} \\ \hline
        \textbf{№} & \textbf{Назва поля} & \textbf{Тип даних} \\ \hline
        \endhead

        1 & id & INTEGER (PK) \\ \hline
        2 & client\_name & TEXT \\ \hline
        3 & car\_info & TEXT \\ \hline
        4 & labor\_cost & REAL (FLOAT) \\ \hline
        5 & parts\_cost & REAL (FLOAT) \\ \hline
        6 & status & VARCHAR \\ \hline
    \end{xltabular}
\end{spacing}

\customSubSection[sec:analysis_third_level_section]{Міждокументні посилання}{1em}

Як зазначено в підрозділі \ref{sec:analysis_two_level_section}, аналіз предметної області та вибір стеку технологій (Python/FastAPI) є критично важливими для швидкодії системи.

На \textbf{рис.~\ref{fig:dashboard_view}} продемонстровано реалізацію веб-інтерфейсу, що взаємодіє з сервером через асинхронні запити.

Використання FastAPI для побудови сучасних API детально розглядається у джерелах \citecustom{1}, \citecustom{3}.

Дані про структуру таблиць бази даних наведені у \tabref{tab:db_orders}.

\begin{spacing}{1.3}
    \noindent
    \begin{xltabular}{\textwidth}{| c | >{\raggedright\arraybackslash}X | c |}
        \caption{Перелік статусів замовлень} \label{tab:statuses} \\
        \hline
        \textbf{№} & \textbf{Статус} & \textbf{Опис} \\ \hline
        \endfirsthead

        \multicolumn{3}{r}{\textit{}} \\ \hline
        \textbf{№} & \textbf{Статус} & \textbf{Опис} \\ \hline
        \endhead

        1 & В роботі & Замовлення щойно створене або виконується \\ \hline
        2 & Очікує запчастин & Необхідна закупівля деталей \\ \hline
        3 & Готово & Ремонт завершено \\ \hline
        4 & Видано & Автомобіль повернуто клієнту \\ \hline
    \end{xltabular}
\end{spacing}
